% TOP_PROBLEM: {"id": 30006797, "title": "Wall's Poincare Duality Group Conjecture", "category_id": 17, "category": {"id": 17, "name": "group_theory", "display_name": "Group Theory", "description": "Problems about groups, group actions, representations, and related algebraic structures.", "slug": "group-theory", "order_index": 17, "created_at": "2026-07-31T15:26:25.671Z"}, "status": "open"}
% FIRST_POSTED: 2026-09-27
% !TeX program = lualatex
% Compile twice: lualatex 30006797.tex
% All references and prose are embedded; no BibTeX database or external inputs.
\documentclass[11pt,a4paper]{article}
\usepackage[margin=24mm,headheight=14pt]{geometry}
\usepackage{fontspec}
\setmainfont{Latin Modern Roman}
\setsansfont{Latin Modern Sans}
\setmonofont{Latin Modern Mono}
\usepackage{amsmath,amssymb,mathtools}
\usepackage{unicode-math}
\setmathfont{Latin Modern Math}
\usepackage{microtype}
\usepackage{enumitem,xurl,needspace}
\usepackage[dvipsnames]{xcolor}
\usepackage{hyperref}
\hypersetup{unicode=true,colorlinks=true,linkcolor=MidnightBlue,urlcolor=MidnightBlue,
 pdftitle={Research Notebook: Section 30006797},pdfauthor={Research notebook},bookmarksnumbered=true}
\usepackage{bookmark}
\usepackage{tocloft}
\setlength{\cftsecnumwidth}{3em}
\cftsetpnumwidth{2.5em}
\cftsetrmarg{4em}
\usepackage{titlesec}
\titleformat{\section}{\Large\bfseries\raggedright}{\thesection}{0.7em}{}
\usepackage{fancyhdr}
\pagestyle{fancy}\fancyhf{}
\fancyhead[L]{\small\sffamily Open Problems | Research notebook}
\fancyhead[R]{\small 221 | 27 September 2026}
\fancyfoot[C]{\thepage}
\setlength{\parindent}{0pt}
\setlength{\parskip}{5pt plus 1pt}
\setlength{\emergencystretch}{3em}
\setcounter{tocdepth}{1}
\setcounter{secnumdepth}{1}
\setlist[itemize]{leftmargin=3em,itemsep=5pt,topsep=4pt,parsep=0pt}
\allowdisplaybreaks
\newcommand{\N}{\mathbb{N}}
\newcommand{\Z}{\mathbb{Z}}
\newcommand{\Q}{\mathbb{Q}}
\newcommand{\R}{\mathbb{R}}
\newcommand{\C}{\mathbb{C}}
\newcommand{\F}{\mathbb{F}}
\newcommand{\A}{\mathbb{A}}
\newcommand{\PP}{\mathbb P}
\newcommand{\E}{\mathbb{E}}
\newcommand{\eps}{\varepsilon}
\newcommand{\dd}{\,\mathrm d}
\newcommand{\sref}[2]{\hyperlink{source:#1:#2}{[S#2]}}
\newcommand{\eref}[2]{\hyperlink{extra:#1:#2}{[E#2]}}
% Turn the next switch off for a shorter reading copy. The quotations preserve
% every exactTarget field, including qualifications omitted from a short summary.
\newif\ifshowcatalogue
\showcataloguetrue
\newcommand{\cataloguescope}[1]{\ifshowcatalogue
 \paragraph{Exact scope in the supplied catalogue.}
 \begin{quote}\small #1\end{quote}\fi}

% Independently compilable extraction; original section text retained.
\numberwithin{equation}{section}
\begin{document}
\setcounter{section}{0}
% ================================================================
% problemId: 30006797
% Canonical problem ID: 30006797
% exactTarget SHA-256: 8ca533b6e8f120d31cf6508d4acec3c1bc3690b40ba06f89165a54bb8fdb595d
\clearpage
\section*{\texorpdfstring{Wall's Poincare Duality Group Conjecture}{Wall's Poincare Duality Group Conjecture}}\label{30006797}
\subsection{Definitions and mathematical statement}
An integral Poincar\'e duality group of dimension $n$ is a group of cohomological dimension $n$, of type FP over $\Z$, whose cohomology with group-ring coefficients vanishes outside degree $n$ and is infinite cyclic in degree $n$. The action on that cyclic dualizing module defines its orientation character $w:G\to\{\pm1\}$. For every finitely presented such $G$ and every $n\ge3$, ask for a closed aspherical topological $n$-manifold $M$ with
\[
 \pi_1(M)\cong G,
\]
inducing the same orientation character. Aspherical means the universal cover is contractible. A finite Poincar\'e duality complex is not automatically a topological manifold, and that distinction is central to the target.
\par\smallskip\textit{Formulation and concept references:} \sref{30006797}{1}.
\cataloguescope{For every n\ensuremath{\ge}3 and every finitely presented integral Poincaré duality group G of dimension n, prove that there is a closed aspherical topological n-manifold M with \ensuremath{\pi}1(M)\ensuremath{\cong}G and matching orientation character.}
\subsection{Short English statement}
Is every finitely presented integral Poincaré duality group of dimension at least three the fundamental group of a closed aspherical manifold of that dimension?
% BEGIN RESEARCH ATTEMPT 221
\subsection{Research attempt}

\paragraph{Current frontier.}
A finitely presented integral Poincar\'e duality group admits a finitely
dominated aspherical Poincar\'e complex. Passing to a finite complex has a
$\widetilde K_0(\Z G)$ obstruction; passing from a Poincar\'e complex to a
manifold has further surgery obstructions. These steps are distinguished
in L\"uck's Theorem~5.4, Lemma~5.7 and Remark~5.8.
\sref{30006797}{1}. The June 2026 revision of Davis--Hillman classifies
compact aspherical four-manifolds with boundary and elementary amenable
fundamental group. Its restricted classification does not realize every
abstract finitely presented PD group. \eref{30006797}{1}, Abstract.
The present attempt keeps all dimensions $n\ge3$ and the orientation
character, and tests whether duality itself could annihilate the first
obstruction.

\paragraph{Attack route.}
Three approaches suggest themselves. A finite projective resolution might
be converted to a finite CW model; this needs vanishing of its reduced
projective class, not just duality of the resolution. Stabilizing by a
circle kills that class, but then one must \emph{destabilize} the resulting
space or manifold. Finally, surgery can act on a finite Poincar\'e complex,
but requires a normal reduction and the correct vanishing obstruction.
We follow the first two routes explicitly and use the third as a precise
conditional completion, not as a theorem that the missing inputs exist.

\paragraph{Attempt.}
\textbf{1. The exact restriction imposed by duality.}
Let $w:G\to\{\pm1\}$ be the orientation character and let
$C_*$ be a bounded finite-projective $\Z G$ chain model for a finitely
dominated $BG$. The construction of such a Poincar\'e model is the input
from \sref{30006797}{1}, Theorem~5.4. The involution
\[
 \overline g=w(g)g^{-1}
\]
turns the dual of a left projective module back into a left module; write
$[P]^*=[\operatorname{Hom}_{\Z G}(P,\Z G)]$ with this convention.
The finiteness obstruction is
\[
 o(G)=\sum_i(-1)^i[C_i]\in\widetilde K_0(\Z G).
\]
Chain Poincar\'e duality compares $C_*$ with the degree-reversed twisted
dual. Taking alternating projective classes gives the rigorous identity
\begin{equation}\label{30006797:selfdual}
 o(G)=(-1)^n o(G)^*.
\end{equation}
Here a chain homotopy equivalence preserves the alternating class because
its contractible mapping cone has zero alternating class; contractible
bounded projective complexes split into elementary adjacent pairs.

Equation \eqref{30006797:selfdual} is not a vanishing theorem. As an abstract
abelian-group calculation, for even $n$ the group $\Z$ with identity
involution has every element satisfying it. For odd $n$, $\Z/2$ with
identity involution has a nonzero solution. These are counterexamples to
an \emph{algebraic inference}, not examples of group-ring obstructions
realized by PD groups. More generally, decomposing rationally into the
$\pm1$ eigenspaces of $*$ still leaves the permitted eigenspace, and
rationalizing loses possible integral torsion. Neither maneuver proves
$o(G)=0$ in the actual reduced $K$-group.

\textbf{2. Circle stabilization hides exactly this information.}
Let $R=\Z G$ and $R'=R[t,t^{-1}]$. After extending scalars, the chain
complex of $BG\times S^1$ is represented by the mapping cone of
\[
 t-1:C_*\otimes_RR'\longrightarrow C_*\otimes_RR'.
\]
In degree $i$ it has projective module
$(C_i\otimes_RR')\oplus(C_{i-1}\otimes_RR')$. Therefore
\begin{equation}\label{30006797:circle}
 o(BG\times S^1)=\sum_i(-1)^i([C_i\otimes R']+[C_{i-1}\otimes R'])=0.
\end{equation}
This computation does not require $o(G)=0$. By the finiteness-obstruction
criterion recalled in \sref{30006797}{1}, Section~5, the stabilized finitely
dominated space has a finite CW model. It is an aspherical PD space of
dimension $n+1$ and fundamental group $G\times\Z$.

The cancellation is useful but dangerous. It gives no way to recover the
original class from the stabilized class: the map on this obstruction is
identically zero. Even a manifold model for $G\times\Z$ would not by
itself produce a manifold model for $G$ by taking an arbitrary map to
$S^1$ and declaring its fiber a closed aspherical manifold. A continuous
map need not be a bundle, its fibers need not be manifolds, and a
regular-value fiber in a smooth representative need not be aspherical.
A separate fibering or splitting theorem would be necessary.

\textbf{3. A conditional high-dimensional completion.}
Fix $n\ge5$. Suppose the following additional facts can be supplied for
the given $G$: the obstruction $o(G)$ vanishes; the resulting finite
aspherical PD complex $X$ admits a degree-one topological normal map
$f:M\to X$ with the specified orientation character; and the associated
Wall obstruction in the appropriate $L_n(\Z G,w)$ group is zero.
The high-dimensional surgery theorem then modifies $f$, within its
normal bordism class, to a homotopy equivalence $M'\to X$.
Consequently $M'$ is closed, has fundamental group $G$, and its universal
cover is homotopy equivalent to the contractible universal cover of $X$.
It is therefore aspherical and realizes the target, including $w$.
This is an application of the normal-map/surgery framework, not a
construction of the normal map. \eref{30006797}{2}, surgery obstruction and
existence theorem.

There are two different missing inputs even after finiteness. A spherical
Spivak normal fibration need not already have the required stable
\emph{topological bundle} reduction. A normal map, once obtained, may
have a nonzero obstruction. A periodic $L$-theory assembly statement
cannot simply be identified with every connective surgery condition:
L\"uck's Remark~5.8 explicitly records the missing degree-zero issue.
His Lemma~5.7 produces homology ANR manifolds under additional assumptions
and in its stated dimensions; such a space is not automatically the
closed topological manifold requested here. \sref{30006797}{1}.

\textbf{4. Testing simpler numerical invariants.}
In an oriented odd-dimensional finite PD complex, rational duality gives
$b_i=b_{n-i}$ and hence $\chi(X)=0$. This agrees with the manifold
identity but follows without a manifold structure. In even dimension,
duality supplies the symmetric or skew form on the middle-dimensional
homology, but says nothing by itself about a local Euclidean structure.
Thus checking Euler characteristic or a self-dual finite resolution
cannot replace either obstruction calculation above.

An attempted minimal-counterexample argument also stops at a specific
point: passing to $G\times\Z$ raises dimension and erases the finiteness
obstruction rather than decreasing a complexity parameter. It provides
no induction whose hypothesis applies to $G$, and no monotonicity of
manifold realizability under removing a direct $\Z$ factor was proved.

\paragraph{Stress test.}
The toy groups $\Z$ and $\Z/2$ in part~1 are not asserted to occur as
$\widetilde K_0(\Z G)$ for a counterexample PD group. They only invalidate
the inference from an involution identity to zero. Conversely,
\eqref{30006797:circle} is an exact projective-class cancellation, not evidence
that $G$ itself has a finite classifying complex. Orientation twists are
retained in both the duality and surgery expressions.

The surgery completion was explicitly restricted to $n\ge5$.
Dimensions three and four in the original target require additional
arguments and are not swept into high-dimensional surgery. Likewise,
the recent elementary-amenable four-manifold classification presupposes
and restricts the geometric objects under study. No claim is made that
it proves the unrestricted existence assertion.

\paragraph{Outcome.}
\textbf{Outcome: Exploratory approach with unresolved gap.}

The attempt derives the exact duality constraint on the finiteness
obstruction and shows algebraically why circle stabilization cannot
recover its vanishing. It gives a conditional high-dimensional
completion with each geometric input named separately. These are
obstruction calculations, not new realizations or claims of priority.

The missing steps are vanishing of the actual finiteness obstruction,
construction of the appropriate normal data, removal of the surgery
obstruction, and separate treatment of dimensions three and four.
No counterexample PD group, and no proof for all such groups, is supplied.
\needspace{20\baselineskip}
% END RESEARCH ATTEMPT 221
\subsection{Sources}
\begin{itemize}\raggedright
\item[\textnormal{[S1]}]\hypertarget{source:30006797:1}{} Wolfgang Lück, “Survey on Aspherical Manifolds,” arXiv:0902.2480v3 (2009); published in European Mathematical Society proceedings (2010).
\url{https://arxiv.org/pdf/0902.2480}.
{\small\textit{Catalogue formulation source.}}
% BEGIN ADDED REFERENCES 221
\item[\textnormal{[E1]}]\hypertarget{extra:30006797:1}{}
J. F. Davis and J. A. Hillman, \emph{Aspherical 4-manifolds with elementary
amenable fundamental group}, arXiv:2501.12512v2, 29 June 2026, Abstract
and Introduction. \url{https://arxiv.org/abs/2501.12512v2}.
Consulted 27 September 2026; its boundary and group-class restrictions
are not omitted here.
\item[\textnormal{[E2]}]\hypertarget{extra:30006797:2}{}
C. T. C. Wall, \emph{Surgery on Compact Manifolds}, 2nd ed., edited by
A. A. Ranicki, Mathematical Surveys and Monographs 69, AMS, 1999,
the normal-map obstruction and high-dimensional surgery existence
chapters. \url{https://www.maths.ed.ac.uk/~v1ranick/books/scm.pdf}.
The application is conditional on the normal data and vanishing obstruction.
% END ADDED REFERENCES 221
\end{itemize}

\end{document}
